% Gerado por tools/tei2latex.py a partir de vol1-divisao-do-ente.xml. Regenerar, não editar.
% Recorte manual: tirada a matéria preliminar que tools/trecho.py arrastou; \pg{331} reposto.

\pg{331}

%%>>> cap-10

\chapter*{CAPITULO SEGUNDO}
\addcontentsline{toc}{chapter}{CAPITULO SEGUNDO}

\textbf{DIVISÃO DO ENTE} \textsc{Summario}: --- Ente necessario e ente contingente. --- Este infinito e ente finito.

--- Ente simples e ente composto. --- Ente immutavel e ente mutavel.

%%>>> art-12

\section*{Artigo I\\ Ente necessario e ente contingente.}
\addcontentsline{toc}{section}{Artigo I. --- Ente necessario e ente contingente.}

\textbf{54. Divisão do ente.} --- Todo o ente --- ou é \emph{acto puro}, --- ou consta de \emph{acto} e de \emph{potencia}. Se é \emph{acto puro}, o ente é \emph{necessario}, \emph{infinito}, \emph{simples}, \emph{immutavel}. Se consta de \emph{acto} e de \emph{potencia}, é \emph{contingente}, \emph{finito}, \emph{composto}, \emph{mutavel}. D'onde a divisão de todo o ente --- em \emph{necessario} e \emph{contingente}, --- em \emph{infinito} e \emph{finito}, --- em \emph{simples} e \emph{composto}, --- em \emph{immutavel} e \emph{mutavel}. --- Esta divisão funda-se na diversa perfeição do ente (1). ---\footnote{Todo o ente é dotado de \emph{perfeição}. --- \emph{Perfeição é toda a realidade}, que \emph{serve para} fazer ou completar uma coisa... A perfeição pode ser \emph{parcial} e \emph{total}. \emph{É parcial}, se representa \emph{uma} das realidades, de que o ente é dotado. \emph{É total}, se representa o \emph{conjuncto} das mesmas realidades. D'onde se vê que a perfeição não é um attributo especial do ente, mas é uma coisa que constitue, adequada ou inadequadamente, o proprio ente. --- Por isso um ente é e diz-se \emph{perfeito}, isto é, \emph{inteiramente feito}, emquanto possue a perfeição \emph{total}, de modo que não lhe falte nenhuma das realidades necessarias para o seu complemento. --- A' \emph{perfeição} oppõe-se a \emph{imperfeição}, que denota \emph{ausencia} de perfeição. A imperfeição é \emph{privativa} ou \emph{negativa}. É \emph{privativa}, quando é ausencia de uma perfeição \emph{devida}, e diz-se commummente \emph{privação}; tal é a ausencia de um braço no homem. É \emph{negativa}, quando é a ausencia de uma perfeição \emph{não-devida}, e chama-se \emph{negação}; tal é a ausencia das azas no homem, ou da razão na planta.\par A perfeição \emph{total} divide-se em \emph{absoluta} e \emph{relativa}. \emph{A absoluta} é o \emph{conjuncto} de todas as realidades ou perfeições \emph{possiveis}, sem limitação alguma; tal é a perfeição de Deus, que é o proprio Ser subsistente. \emph{A relativa} é o conjuncto das realidades ou perfeições, que são devidas a um certo genero, ou a uma certa especie de entes: tal é a perfeição de um homem, que tem tudo o que a sua natureza exige. --- A perfeição \emph{absoluta} não admitte graus, nem divisões: a \emph{relativa} admitte graus, --- emquanto um genero ou uma especie pode abranger um maior ou menor numero de perfeições parciaes com relação a um outro genero, ou a uma outra especie, --- ou emquanto um individuo possue uma realidade ou perfeição por um modo mais \emph{intenso} que o \pgn{332}modo por que a possue um outro individuo do mesmo genero, ou da mesma especie. --- A perfeição \emph{total} pode considerar-se --- em relação á \emph{entidade}, quando o ente tem tudo o que é indispensavel para a sua essencia e integridade; --- em relação á \emph{operação}, quando o ente tem tudo o que é necessario para operar em conformidade com a sua essencia; --- em relação ao \emph{fim}, quando o ente, tendo alcançado o fim pelas suas operações, encontra nelle o seu complemento, pois o fim é bem (Cf. \emph{Sum. Th.}, p. I, q. 6, a. 3).\par A perfeição \emph{relativa} divide-se em \emph{simples} e \emph{mista}. --- \emph{Simples} é a que no seu conceito não encerra imperfeição alguma e não se oppõe a perfeições mais elevadas; tal é a \emph{vida}, a \emph{bondade}, a \emph{belleza}, etc. A perfeição \emph{simples} costuma chamar-se tambem \emph{absoluta}, mas emquanto o seu conceito não denota imperfeição. --- \emph{Mista} é a que, embora no seu conceito exprima uma entidade positiva, todavia encerra alguma imperfeição e se oppõe a uma perfeição mais elevada: tal é a \emph{corporeidade}, que, embora denote uma realidade, comtudo importa uma imperfeição, pois o \emph{corpo} é \emph{composto}, \emph{dependente}, etc., --- e se oppõe á \emph{espiritualidade}, que é uma coisa mais elevada do que o corpo (Cf. Ansel., Monol., c. 15).\par A perfeição pode convir ao ente por tres modos: \emph{formal}, \emph{virtual} e \emph{eminente}. Convem \emph{formalmente}, quando existe no ente segundo a propria forma ou razão \emph{especifica}, isto é, segundo o seu conceito e a sua definição; assim encontra-se o \emph{calor no fogo}, a \emph{sensibilidade no cão}, a \emph{razão no homem}. Convem \emph{virtualmente}, quando não existe no ente segundo a propria forma, mas existe na virtude ou causalidade do ente, emquanto este é capaz de a produzir; assim encontra-se a \emph{planta na semente}, e todo o \emph{effeito na sua causa}. Convem \emph{eminentemente}; quando existe no ente, nem segundo a propria forma, nem segundo a energia ou causalidade, mas por um modo mais nobre e mais elevado que o modo por que se encontra nos seres, que a possuem \emph{formalmente}; assim encontram-se em Deus todas as perfeições dos entes creados. --- Note-se que um d'esses modos não exclue outro. Uma perfeição pode convir a um ente por um modo \emph{formal} e \emph{eminente} ao mesmo tempo, e convem quando existe nelle segundo a propria forma, mas num grau mais elevado; assim a \emph{sabedoria}, que é o conhecimento das coisas pelas causas supremas, attribue-se a Deus --- não só \emph{formalmente}, emquanto, em Deus, como no homem, importa o conhecimento pelas causas supremas, --- mas tambem \emph{eminentemente}, emquanto em Deus a sabedoria é uma e a mesma coisa com essencia, ao passo que no homem é uma qualidade accidental e limitada. Como tambem uma perfeição pode convir a um ente por um modo \emph{virtual} e \emph{eminente} ao mesmo tempo, e convem quando não existe nelle segundo a propria forma, mas existe na virtude ou capacidade do ente, e de um modo mais elevado e nobre; assim \emph{o raciocinio}, importando um discurso ou uma passagem da potencia para o acto, não se attribue a Deus \emph{formalmente}, --- mas attribue-se \emph{virtualmente}, emquanto a intelligencia divina tem a capacidade de perceber a conclusão no principio, --- e tambem \emph{eminentemente}, emquanto a intelligencia\pgn{333} divina não passa, como a do homem, do conhecido para o desconhecido, mas vê a conclusão no principio sem discurso, por intuição.\par Premissas estas noções, vejamos a razão da divisão do \textbf{ente}. A perfeição pode encontrar-se no ente --- ou de um modo \emph{necessario} ou \emph{contingente}, e temos o ente \emph{necessario} e o ente \emph{contingente}, --- ou de um modo \emph{illimitado} ou \emph{limitado}, e temos o ente \emph{infinito} e o ente \emph{finito}, --- ou de um modo \emph{simples} ou \emph{composto}, e temos o ente \emph{simples} e o ente \emph{composto}, --- ou de um modo \emph{permanente} ou \emph{variavel}, e temos o ente \emph{immutavel} e o ente \emph{mutavel}. --- A primeira especie da divisão refere-se á \emph{existencia} do ente; a segunda, á sua \emph{comprehensão}; a terceira, á sua \emph{constituição}; a quarta, á \emph{permanencia} ou \emph{estabilidade} da perfeição. --- Como se vê, todas estas especies denotam a mesma coisa e importam os mesmos extremos, considerados sob differentes aspectos e significados por varios conceitos e termos, dos quaes um é mais claro e mais obvio que outro: porque o Ente \emph{necessario}, \emph{infinito}, \emph{simples}, \emph{immutavel} é Deus; e o ente \emph{contingente}, \emph{finito}, \emph{composto}, \emph{mutavel} é a creatura. Por isso esta divisão é a mais transcendental que se pode dar. --- Advertimos que, tratando aqui do Ente \emph{necessario}, \emph{infinito}, \emph{simples} e \emph{immutavel}, não entendemos demonstrar a sua existencia; mas, admittida, pela experiencia, a existencia de entes \emph{contingentes}, \emph{finitos}, \emph{compostos} e \emph{mutaveis}, procuramos descobrir o nexo, que elles têm com o Ente \emph{necessario}, \emph{infinito}, \emph{simples} e \emph{immutavel}.}
\pg{332} 55. \textbf{Ente necessario e ente contingente}. --- Ente necessario é o que existe \emph{por si}, \emph{pela exigencia da sua essencia}, e por isso não pode deixar de existir. --- Ente \emph{contingente} é o que existe \emph{accidentalmente}, não pela exigencia da sua essencia, mas pela causalidade de outro ente, e por isso é indifferente para a\pg{333} existencia e para a não-existencia. --- No ente \emph{necessario}, a \emph{existencia} é um elemento da \emph{essencia}, e identifica-se com esta; no ente \emph{contingente}, a \emph{existencia} é \emph{indifferente} para a \emph{essencia}, e uma coisa é \emph{realmente} distincta de outra \footnote{\emph{Ente necessario} é o que não pode deixar de ser. Na verdade, \emph{ente} é o que possue o \emph{ser}, é o que tem a existencia, é o \emph{existente}. Portanto, o ente \emph{necessario} é o que existe \emph{necessariamente}, e por isso possue o \emph{ser} por tal modo, que lhe repugna o \emph{não-ser}. --- O \emph{ente contingente}, por ser \emph{contingente}, existe por tal modo, que pode deixar de existir, sem que d'ahi resulte algum absurdo ou inconveniente. --- O ente \emph{necessario} pode ser \emph{absoluto} e \emph{hypothetico}. \emph{É absoluto}, quando o \emph{não-ser} lhe repugna pela sua intrinseca natureza, independentemente de toda a condição ou hypothese; tal é Deus. É \emph{hypothetico}, quando o \emph{não-ser} lhe repugna, não pela sua intrinseca natureza, mas debaixo de uma certa condição ou hypothese; tal é toda a creatura, emquanto existe. --- Aqui fallamos do ente \emph{absoluto}.\par O ente absoluto, não tendo nenhuma potencia ou indifferença para a existencia, mas existindo em virtude da propria essencia, é ente \emph{por si, pela sua essencia}, de modo que a sua essencia é a razão sufficiente da sua existencia. --- Note-se que quando se diz que o ente \emph{absoluto} existe pela \emph{sua essencia}, não se entende que a essencia d'esse ente é a \emph{causa} da sua existencia: o que é absurdo; mas deve entender-se que esse ente não depende, na sua existencia, de uma causa extrinseca e superior. --- Advirta-se tambem que, embora a \emph{essencia} do ente \emph{absoluto} se identifique com a \emph{existencia}, todavia a nossa intelligencia não percebe essas duas coisas por uma só idea, mas por duas ideas distinctas, uma das quaes representa a \emph{essencia}, e a outra a \emph{existencia}. E nada mais natural. A intelligencia humana não tem um conhecimento \emph{intuitivo do} Ente \emph{absoluto}, mas eleva-se a Elle por meio das coisas creadas, e neslas a \emph{essencia} é \emph{realmente} distincta da \emph{existencia}.}.
\pg{334}
\textbf{56. O ente, que é \emph{acto puro}, é \emph{necessario}; e o ente, que é \emph{necessario}, é \emph{acto puro}.}

\textbf{a)} \emph{O ente, que é acto puro, é necessario.} --- O ente, que é \emph{acto puro}, exclue de si toda a potencia; aliás não seria o que é. Ora o ente, que exclue de si toda a potencia, deve possuir a existencia identificada com a essencia, e por isso é \emph{necessario}. Logo o ente, que é \emph{acto puro}, é \emph{necessario}\footnote{O \emph{acto puro} exclue de si toda a potencia. Ora o ente, que exclue de si toda a potencia, não pode ser \emph{contingente}; porque o ente \emph{contingente} não exclue, mas encerra a potencia, emquanto está em potencia para o acto da existencia. Portanto o ente, que é \emph{acto puro}, é \emph{necessario}, possue a existencia pela exigencia da sua essencia. Nem podia ser diversamente. O ente, que é \emph{acto puro}, possue a existencia identificada com a essencia, de modo que podemos dizer que a existencia é um dos elementos da essencia. Sendo assim, esse ente é \emph{necessariamente existente}, como o homem, que pela sua essencia é dotado de razão, é \emph{necessariamente racional}. As essencias, como vimos, são \emph{necessarias} quanto aos elementos, que as constituem.}.

\textbf{b)} \emph{O ente, que é necessario, é acto puro.} --- O ente, que é \emph{necessario}, existe pela exigencia da sua essencia, e por isso possue a existencia identificada com a essencia. Ora o ente, que possue a existencia identificada com a essencia, é \emph{acto puro}. Logo o ente, que é \emph{necessario}, é \emph{acto puro}\footnote{O ente \emph{necessario} é \emph{acto puro}. O \emph{acto} é puro, quando não é recebido num sujeito distincto, nem está em potencia para um acto ulterior. Ora o ente \emph{necessario} --- não é recebido num sujeito distincto, porque o seu \emph{ser} é uma e a mesma coisa com a \emph{essencia}, --- nem está em potencia para um acto ulterior, porque o ente, cujo ser se identifica com a essencia, é o proprio ser, é o ser subsistente, e por isso possue a plenitude da perfeição.

\par D' onde se vê que o ente \emph{necessario} é \emph{unico}. Por quanto o \emph{acto}, como dissemos, só se multiplica emquanto é recebido em muitos sujeitos. Ora o ser do ente \emph{necessario}, identificando-se com a essencia, não é recebido em nenhum sujeito, mas é por si subsistente, e por isso não se multiplica, mas é \emph{unico}. Logo o ente \emph{necessario} é \emph{unico}.}.

\textbf{57. O ente, que consta de acto e de potencia, é contingente; e o ente, que é contingente, consta de acto e de potencia.}

\textbf{a)} \emph{O ente, que consta de acto e de potencia, é contingente.} --- O ente, que consta de \emph{acto} e de \emph{potencia}, possue a existencia \emph{realmente} distincta da essencia. Ora o ente, que possue a existencia \emph{realmente} distincta da essencia, é \emph{contingente}. Logo o ente, que consta de \emph{acto} e de \emph{potencia}, é \emph{contingente}\footnote{A conclusão é certissima. Na verdade, todo o ente composto pode ser dissolvido; e, podendo ser dissolvido, é naturalmente \emph{contingente}. --- Além d'isso, o ente composto é \emph{posterior} ás partes componentes, e por isso não pode \emph{ser necessario}, mas deve ser \emph{contingente}: porque o ente \emph{necessario},\pgn{335} sendo o proprio ser, não está ordenado para outro ser mais perfeito, e é simplesmente o \emph{primeiro}. --- Mais: o composto depende, no seu ser, de uma causa, que une as parte componentes e que deve ser não só anterior ao proprio composto, mas distincta delle. Dependendo, no seu ser, de uma causa distincta, o composto não existe pela sua natureza, e por isso é \emph{contingente}.}\pg{335}b) \emph{O ente, que é contingente, consta de acto e} \textbf{\emph{de potencia.}} — O ente, que é \emph{contingente}, não existe pela exigencia da sua essencia. Ora o ente, que não existe pela exigencia da sua essencia, consta de existencia—\emph{realmente}—distincta—da essencia, e por isso de \emph{acto} e de \emph{potencia}. Logo o ente, que é \emph{contingente}, consta de \emph{acto} e de \emph{potencia} \footnote{(1) A essencia do ente \emph{contingente}, mesmo quando é actual e existe na realidade, distingue-se realmente da \emph{existencia}, como a potencia real se distingue do acto real. Portanto o ente \emph{contingente} consta de \emph{acto} e de \emph{potencia}. — D'ahi se segue que o ente \emph{contingente} não é necessariamente unico, mas podem existir muitos entes \emph{contingentes}. Por quanto, o ser do \emph{contingente}, sendo recebido na essencia, pode multiplicar-se indefinitamente, conforme a multidão dos sujeitos, ou das essencias, em que é recebido.}.

\textbf{58. O ente necessario é eterno.} — O ente \emph{necessario} é o que, pela exigencia da sua essencia, não pode deixar de existir. Ora o ente, que não pode deixar de existir, existiu sempre, existirá sempre, é \emph{eterno}. Logo o ente \emph{necessario} é \emph{eterno}.

\textbf{59. O ente contingente é produzido por uma causa extrinseca.} — Tudo o que se encontra num sujeito, e não pertence á sua essencia, deve ser produzido — ou pelos principios essenciaes do proprio sujeito — ou por uma causa extrinseca. Não ha meio termo. Ora a existencia, que não pertence á essencia do ente \emph{contingente}, não pode ser produzida pelos principios essenciaes do proprio ente; porque, neste caso, a essencia preexistiria á existencia (pois a causa preexiste ao effeito): o que é absurdo. Logo deve ser produzida por uma causa extrinseca. Logo o ente \emph{contingente} é produzido por uma causa extrinseca\footnote{(2) O ente \emph{necessario}, não podendo deixar de ser, é \emph{eterno}; porque, se não fosse eterno, se tivesse começo ou fim, seria \emph{contingente}. — Além d'isso, o ente \emph{necessario}, sendo \emph{acto purissimo}, não pode ter nem principio, nem fim, nem successão, porque tudo isto envolve potencialidade; e por isso é \emph{eterno}, porque a eternidade exclue principio, fim e successão. — Não só o ente \emph{necessario} é \emph{eterno}, mas é a sua \emph{eternidade}; porque elle é o seu ser, e a sua duração ou permanencia no \emph{ser} não se distingue da sua essencia, e por isso, assim como é a sua essencia, assim tambem é a sua eternidade. (\emph{Sum}. \emph{Th}. p. I, q. 10, a. 2; \emph{c. Gent.}, l. I, c. 15).\par (3) O ente \emph{contingente}, por não ser o proprio ser ou a propria existencia, isto é, por não ser \emph{necessario}, é \emph{indifferente} para a existencia e para a não-\pgn{336}existencia, e por isso não pode ser eterno, e teve começo. — E como recebeu elle a existencia real? Foi pela exigencia da sua natureza? Não. Porquanto, um sujeito, que é por si indifferente para dois attributos oppostos, não pode ter na sua essencia a razão pela qual exija um attributo, e não o opposto. Ora o ente \emph{contingente} é, pela sua essencia, indifferente para a existencia e para a não-existencia, é por isso a existencia não lhe pode convir pela mesma essencia. Logo o ente \emph{contingente} deve receber de um agente extrinseco a existencia real. (\emph{Sum}. \emph{Th}., p. I, q. 3, a. 4).}
\pg{336}

\textbf{60. O ente contingente é produzido, immediata ou mediatamente, pelo ente necessario.}

\textbf{a)} A causa extrinseca, que produz o ente \emph{contingente}, é um

ente \emph{necessario}; ou \emph{contingente}. Se é \emph{necessario}, está provada a conclusão. Se é \emph{contingente}, este foi produzido por um ente \emph{necessario}, ou por um ente \emph{contingente}. Se foi produzido por um ente \emph{necessario}, temos o que queriamos. Se foi produzido por um ente \emph{contingente}, renovamos a pergunta, e, subindo de um ente para outro, devemos inevitavelmente encontrar o ente \emph{necessario}; porque, se não existisse o ente \emph{necessario}, não poderia existir nenhum ente \emph{contingente}, que não tem em si a razão da sua existencia. Logo o ente \emph{contingente} é produzido, immediata ou mediatamente, pelo ente \emph{necessario} \footnote{A conclusão é um corollario da precedente. O ente contingente, por ser \emph{contingente}, não tem em si a razão da propria existencia; logo deve tel-a numa causa extrinseca, no ente \emph{necessario}. --- Dizemos que o ente \emph{contingente} depende, \emph{immediata} ou \emph{mediatamente}, do ente \emph{necessario}; porque, embora um ente \emph{contingente} possa ser produzido immediatamente por um outro ente \emph{contingente}, comtudo a ultima razão adequada da existencia do mesmo ente \emph{contingente} não pode encontrar-se senão no ente \emph{necessario}.}.

\textbf{b)} O ente, que se acha em potencia em relação a um certo acto, não o pode adquirir, se não lh'o fôr communicado, immediata ou mediatamente, por outro ente, que o possua pela sua essencia. Ora o ente \emph{contingente}, sendo por si indifferente para a existencia e para a não existencia, não existiu sempre e não pode adquirir a existencia, se não lh'a fôr communicada, immediata ou mediatamente, pelo ente \emph{necessario}, que a possua pela exigencia da sua essencia. Logo o ente \emph{contingente} é produzido, immediata ou mediatamente, pelo ente \emph{necessario} \footnote{O ente \emph{contingente} é um ente \emph{em potencia}, isto é, \emph{pode} receber o acto da existencia, mas não o possue por si, pela exigencia da sua natureza. Ora um ente, que não possue por si uma perfeição, não a pode communicar a si mesmo, isto é, não pode passar por si da potencia para o acto. Logo essa perfeição deve ser-lhe communicada por outro ente. Este ente não poderia\pgn{337} communicar a outro uma perfeição, se a não possuisse, porque ninguem dá o que não tem. Porquanto, o ente \emph{contingente}, que é indifferente para a existencia e a não-existencia, que está em potencia em relação á existencia, ficaria sempre nessa indifferença, nunca poderia adquirir a existencia, se não lhe fosse communicada pelo ente que a possue pela sua essencia, isto é, pelo ente \emph{necessario}. (Cf. \emph{C. Gent}, I, 16). --- Alem d'isso, o ente \emph{contingente} é ente \emph{por participação}, ou \emph{participado}; porque não possue o ser pela propria essencia, não é o proprio ser. Ora o ente \emph{participado} deve derivar do ente, que é tal pela sua essencia; isto é, do ente \emph{imparticipado}; porque o ente \emph{participado}, não sendo o proprio ser, recebeu o ser de outro ente, que, em ultima analyse, deve possuir o ser pela propria essencia. Diz S. Thomaz: « Necesse est dicere omne id quod praeter \textsc{Deum} est, a \textsc{Deo} factum esse. Solus enim Deus est suum \emph{esse}; in omnibus autem aliis differt essentia rei et esse ejus... Et ex hoc manifestum est quod solus Deus est ens per suam essentiam; omnia vero alia sunt entia per participationem. Omne autem quod est per participationem, causatur ab eo quod est per essentiam; sicut omne ignitum causatur ab igne » (\emph{Sum. Th.}, p. I, q. 61, a. 1; Cf. \emph{C. Gent.}, I, 15).}.
\pg{337} 61. Corollarios. --- Dos principios expostos resulta:

I. O ente \emph{necessario}, assim como na sua essencia \emph{ideal} contem a razão da sua existencia \emph{ideal}, assim tambem na sua essencia \emph{real} contem a razão da sua existencia \emph{real}; porque a existencia pertence a mesma ordem, á qual pertence a \emph{essencia} (\textsuperscript{1}).

II. Acertada a existencia \emph{real} de um só ente \emph{contingente}, deduz-se legitimamente a existencia do ente \emph{necessario}; pois o ente \emph{contingente} não tem em si a razão da sua existencia \footnote{(1) O attributo segue sempre a condição do sujeito e por isso encontra-se na mesma ordem, em que se encontra o sejeito. --- Encontra-se o sujeito na ordem \emph{ideal}? O attributo deve encontrar-se na ordem \emph{ideal}. --- Encontra-se o sujeito na ordem \emph{real}? O attributo tambem deve encontrar-se na ordem \emph{real}. Ao ente \emph{necessario} convem \emph{essencialmente} a \emph{existencia}. Portanto, onde se encontra o ente \emph{necessario}; ahi deve encontrar-se a existencia. Mas se o ente \emph{necessario} existisse \emph{exclusivamente} na intelligencia, não poderia ter senão uma existencia \emph{ideal}. Por isso, quando se quer deduzir da \emph{idea} do ente \emph{necessario} a sua existencia \emph{real}, commette-se um sophisma. --- A existencia pertence \emph{essencialmente} á \emph{idea} de ente \emph{necessario}. Por isso, se esse ente existe, \emph{essencial} e \emph{necessariamente} existe. Mas é necessario provar que effectivamente o ente \emph{necessario} existe. Assim, pertencendo a razão á essencia do homem, segue-se que, se existe o homem, existe uma creatura, a qual é essencial e necessariamente racional; mas do facto de ser o homem essencial e necessariamente racional não segue que essa creatura existe na realidade. Essa existencia deve provar-se com argumentos fundados, não nas ideas, mas nos factos.}.\footnote{(2) A demonstração da existencia do ente \emph{necessario} deve fundar-se nos factos. Existem no mundo entes \emph{contingentes}? A experiencia interna e externa diz que existem; porque ha no mundo entes que começam e acabam, que estão sujeitos a mudanças, e por isso apresentam um caracter de \emph{contingencia}. Se existem os entes \emph{contingentes}, deve existir o ente \emph{necessario}. O \pgn{338}proprio \emph{Kant} escreve: « Desde que existe alguma coisa, não podemos deixar de concluir que existe tambem um ente \emph{necessario} ». --- Nem se diga que não podemos concluir da existencia do ente \emph{contingente} a existencia do ente \emph{necessario}, por não haver nenhuma relação entre o \emph{contingente} e o \emph{necessario}. Entre o \emph{contingente} e o \emph{necessario} existe, como vimos, a relação, que existe entre o effeito e a causa; e por isso, assim como da existencia do effeito deduzimos legitimamente a existencia da causa respectiva, assim tambem da existencia do ente \emph{contingente} deduzimos a existencia do ente \emph{necessario}.}
\pg{338} III. Na ordem \emph{ontologica}, ou da \emph{realidade}, o ente \emph{necessario} é \emph{anterior} ao \emph{contingente}, pois este é produzido por aquelle; mas na ordem \emph{logica}, ou do \emph{conhecimento}, o ente \emph{contingente} é \emph{anterior} ao \emph{necessario}, pois é por meio do \emph{contingente} que subimos ao conhecimento do \emph{necessario} \footnote{Na ordem \emph{ontologica}, o ente \emph{necessario} é anterior ao \emph{contingente}, pois aquelle é causa d'este, e toda a causa é anterior ao effeito. --- Na ordem \emph{logica}, porém, dá-se o contrario. O conceito do ente \emph{necessario} é posterior ao conceito do \emph{contingente}. Porquanto, o nosso conhecimento começa pela percepção das coisas sensiveis, que estão mais ao nosso alcance. A nossa intelligencia, penetrando no ámago das coisas sensiveis, descobre nellas caracteres de mutabilidade, de defectibilidade, de potencialidade, e por isso de indifferença para a existencia e para a não-existencia, e assim forma a idea do ente \emph{contingente}. E como o ente, indifferente para a existencia, não pode ser reduzido ao acto da existencia senão pela acção, immediata ou mediata, do ente, que possua por si, pela sua essencia, esse acto, assim da idea e pela idea do ente \emph{contingente} nos elevamos á idea do ente \emph{necessario}.}.

%%>>> art-13

\section*{Artigo II.\\ Ente infinito e ente finito.}
\addcontentsline{toc}{section}{Artigo II. --- Ente infinito e ente finito.}

62. \textbf{Ente infinito e ente finito.} --- Ente \emph{infinito} é o que é dotado de uma realidade ou perfeição \emph{illimitada}, de modo que é a propria perfeição. --- Ente \emph{finito} é o que é dotado de uma realidade ou perfeição \emph{limitada}, de modo que só possue a perfeição conveniente a um genero ou a uma especie. --- A idea do \emph{infinito} é \emph{negativa} quanto ao \emph{termo}, mas \emph{é positiva} quanto ao \emph{objecto}, pois representa um ente dotado de perfeição sem limite. A idea do \emph{finito}, mesmo considerada em relação ao \emph{objecto}, em parte é \emph{positiva}, emquanto exprime uma \emph{perfeição}, --- e em parte é \emph{negativa}, emquanto exprime uma perfeição \emph{limitada}, isto é, unida com a ausencia de ulterior perfeição\footnote{\emph{Infinito} é o que não tem limites. Por isso, o \emph{infinito} não exclue toda e \emph{qualquer} outra realidade, mas só exclue tudo o que o pode limitar\pgn{339}na sua ordem. --- Um ente pode ser \emph{infinito} --- ou quanto á \emph{potencia}, que não é determinada por nenhum acto, como a \emph{materia prima}, --- ou quanto á \emph{forma}, cuja entidade não tem limite. No primeiro caso, o \emph{infinito} diz-se \emph{privativo}, porque denota uma absoluta \emph{privação} ou ausencia de toda e qualquer perfeição num sujeito, que naturalmente devia ter uma perfeição; no segundo caso, o \emph{infinito} diz-se \emph{negativo}, porque exprime uma perfeição não coarctada por limites. --- O infinito \emph{negativo} pode ser \emph{relativo} ou \emph{absoluto}. É \emph{relativo}, quando é illimitado numa certa e determinada ordem de realidade ou perfeição; tal seria uma linha que na ordem do comprimento não tivesse fim. É \emph{absoluto}, quando é illimitado em toda a ordem de perfeição ou realidade, de modo que encerra todas as perfeições possiveis, e é a propria perfeição. É d'este infinito \emph{absoluto}, chamado simplesmente \emph{infinito}, que entendemos fallar neste artigo. --- D'onde se segue que tres elementos concorrem para a formação da idea do \emph{infinito}: um \emph{positivo}, que é a realidade ou perfeição, e dois \emph{negativos}, que são o \emph{limite} e a \emph{negação} do limite. Porquanto, o infinito, em si mesmo, é eminentemente positivo, porque é a propria perfeição; mas a nossa intelligencia não pode concebel-o senão de um modo positivo-negativo, emquanto o concebemos como uma entidade positiva, da qual afastamos todo o limite, que é já uma negação, e concebemos a negação do limite, isto é, a negação da negação. E não podia ser d'outro modo. O infinito, na sua entidade positiva, escapa ao alcance da nossa intelligencia, e por isso, não podendo conceber o que elle é e não podendo exprimil-o pelo conceito, que formamos d'elle e que não é adequado, recorremos á negação, negamos os limites, negamos a negação, e dizemos que \emph{não é limitado}, que é \emph{infinito}. Nominamos a Deus como o concebemos; porque o termo, que denota um objecto, corresponde ao modo, por que o proprio objecto se concebe. Diz S. Thomaz: «\emph{Secundum quod aliquid a nobis cognosci potest, sic a nobis potest nominari}» (\emph{Sum. Th.}, p. I, q. 13, a. 1). --- Do \emph{infinito} distingue-se o \emph{indefinito}; porque este, embora possa augmentar e diminuir sem fim e por isso não tenha limites \emph{determinados}, comtudo na realidade é sempre finito.\par \emph{Finito é o} que tem limites, é o que não tem uma ulterior perfeição. Por isso, o \emph{finito}, como dizemos no texto, mesmo quanto ao \emph{objecto}, importa --- uma \emph{affirmação}, porque \emph{affirma} uma \emph{perfeição}, --- e uma \emph{negação}, porque \emph{nega} uma perfeição \emph{ulterior}. --- O ente \emph{finito} é tal --- quanto á \emph{essencia}, que, pertencendo a um certo genero ou a uma certa especie, não possue a perfeição dos outros generos ou das outras especies, --- e quanto á \emph{existencia}, a qual, por ser recebida numa \emph{potencia limitada}, como é a essencia do ente finito, não pode deixar de ser um acto limitado.}.

\textbf{A nossa intelligencia possue uma idea distincta do infinito.} --- A nossa intelligencia possuirá uma idea \emph{distincta} do \emph{infinito}, se poder apontar o elemento, pelo qual o infinito se distingue dos seres finitos. Ora a nossa intelligencia aponta effectivamente o elemento, pelo qual o infinito se distingue dos seres finitos e que é \emph{a exclusão de todo o limite}. Logo a nossa intelligencia possue uma idea \emph{distincta} do \emph{infinito}. --- Todavia essa idea, embora \emph{distincta}, \emph{é imperfeita} e \emph{inadequada}; porque \pg{339} \pg{340}

uma faculdade finita, como é a nossa intelligencia, não pode medir e comprehender o infinito (1).

\textbf{64. O infinito é cognoscivel e conhece-se effectivamente pelo finito.}

a) O \emph{infinito é cognoscivel pelo finito}. --- No ente finito, ou limitado, devem distinguir-se duas coisas: a \emph{entidade}, ou a perfeição, e a \emph{limitação} d'essa entidade, ou a negação de ulterior entidade. A \emph{entidade limitada} não é limitada emquanto é \emph{entidade}, mas emquanto é \emph{limitada}, isto é, emquanto importa a negação de uma entidade ulterior. Ora a nossa intelligencia, assim como abstrahe de limites determinados e concebe o indefinito, assim tambem pode, sem repugnancia, abstrahir totalmente dos limites, e conceber uma entidade sem limites, isto é, uma entidade infinita. Logo o \emph{infinito} é cognoscivel pelo \emph{finito} \footnote{(1) Que a nossa intelligencia possua a idea do \emph{infinito}, está fora de toda a duvida. Fallamos no infinito, discutimos acerca do infinito; e ninguem falla numa coisa, ninguem discute ácerca de uma coisa, de que não têm a idea. A questão refere-se á \emph{qualidade} d'essa idea. Affirmamos que essa idea é \emph{distincta}, ainda que concedamos que é \emph{imperfeita} e \emph{inadequada}.\par Os \emph{positivistas} negam-nos a possibilidade de conhecer o \emph{infinito}, dizendo que a nossa intelligencia só pode conhecer o finito, e que por isso tudo o que estiver além do finito é por nós incognoscivel. --- O erro dos \emph{positivistas} é manifesto. Se o nosso conhecimento devesse --- ou ser sempre e necessariamente \emph{intuitivo}, --- ou comprehender sempre o seu objecto na sua totalidade, --- ou produzil-o, então a intelligencia não poderia conceber o infinito; porque, sendo ella finita, não pode ter uma intuição directa do infinito, nem comprehendel-o, e muito menos produzil-o. Mas é falso que o nosso conhecimento se refira sempre aos objectos, de que se tenha tido uma percepção intuitiva, como é falso que deva sempre comprehender o seu objecto na sua totalidade, ou que deva produzil-o para o conhecer. Logo a intelligencia pode conhecer o infinito, ainda que dá um modo finito ou inadequado.}.\footnote{(2) São tres, como dissemos, os elementos, que constituem a idea do \emph{infinito}: a \emph{entidade}, o \emph{limite} e a \emph{negação} do limite. Ora a nossa intelligencia, considerando os entes finitos, pode conceber o que é \emph{entidade}, o que é \emph{limite} e o que é \emph{negação} de todo o limite. Pode conceber o que é \emph{entidade}; porque vê que os entes finitos são alguma coisa e não são nada. Pode conceber o que é \emph{limite}; porque vê que um ente é distincto de outro, um é inferior a outro, e assim percebe que nenhum possue a plenitude da entidade, e que todos são limitados a um genero ou a uma especie. Pode conceber a \emph{negação} de todo o limite; porque, assim como, pela sua força de abstracção, pode afastar de um ente uma determinada limitação, ou imperfeição, percebida noutro ente, assim pode afastar de uma entidade todo o limite, toda a imperfeição, e assim conceber, de algum modo, o ente infinito.\par \emph{Descartes}, \emph{Fabre} e todos os sequazes do \emph{ontologismo} dizem que não podemos subir do finito para o infinito, porque este não está contido naquelle.\pgn{341}--- Mas tal raciocinio é vicioso. Para que um ente seja conhecido por outro ente, não é necessario que o primeiro esteja contido no segundo \emph{physica} e \emph{realmente}, mas basta que esteja contido \emph{logicamente}; basta que o segundo tenha com o primeiro uma relação de semelhança ou de dependencia. Ora o finito tem com o infinito uma relação essencial de dependencia e de semelhança; pois, como veremos, o finito é produzido exclusivamente pelo infinito, e reproduz em si, ainda que de um modo immensamente imperfeito, as perfeições do infinito. Do facto de o infinito não estar \emph{physica} ou \emph{realmente} contido no finito só se deduz que o finito não pode dar-nos uma idea \emph{adequada} do infinito, o que concedemos e sustentamos; mas não se deduz que o finito não possa dar-nos nenhuma idea, nem sequer \emph{inadequada}, do infinito; porque, para isso, basta afastar da entidade ou perfeição finita todo e qualquer limite. Será um conhecimento imperfeito do ente perfeitissimo; mas será sempre conhecimento.}
\pg{341} b) O \emph{infinito conhece-se effectivamente} pelo \emph{finito}. --- É certo que temos uma idea distincta do \emph{infinito}, a qual é \emph{positivo-negativa}, emquanto representa uma entidade não-limitada. Essa idea, quanto á parte \emph{negativa}, deriva --- ou \emph{do proprio infinito}, --- ou \emph{da nossa intelligencia}, --- ou \emph{do meio objectivo do conhecimento}. Ora não pode derivar --- nem \emph{do proprio infinito}, pois este é eminentemente positivo, --- nem \emph{da nossa intelligencia}, pois quanto mais positivo é o objecto cognoscivel, tanto mais positivo é o modo por que o conhecemos. Logo deve derivar do \emph{meio objectivo}, de que a intelligencia parte para chegar ao conhecimento do \emph{infinito} e que é o \emph{finito}, emquanto se exclue da sua entidade todo o limite. Logo o \emph{infinito} conhece-se effectivamente pelo \emph{finito} (').\footnote{Este modo de conhecer o \emph{infinito} é muito natural. O primeiro objecto do nosso conhecimento deve ser o que é mais proporcionado com a indole da nossa intelligencia. Ora a nossa intelligencia, sendo finita e limitada, deve primeiramente conhecer os objectos, que são finitos e limitados ; aliás não haveria proporção entre a mesma faculdade e o seu objecto connatural. Logo é do conhecimento do finito que devemos elevar-nos ao conhecimento do infinito. --- Alem d'isso, todo o nosso conhecimento intellectual é precedido pelo conhecimento sensitivo e deriva d'elle, emquanto nada chega á intelligencia, que não tenha passado pelos sentidos. Ora o objecto do conhecimento sensitivo é o ente visivel e finito. Logo o conhecimento do infinito deve ser precedido pelo conhecimento do finito, e deriva d' este.\par Portanto a nossa intelligencia forma a idea do infinito, emquanto exclue da entidade finita todos os limites, porque esta exclusão constitue a essencia do infinito e o extrema dos seres finitos. \emph{Spencer} escreve: «A idea que temos do limitado é composta da idea de uma determinada especie de ser, e da idea dos limites, em que o ser é conhecido. Na sua antithese, na concepção do illimitado, supprime-se a idea dos limites, mas permanece a idea do ser» (\emph{Premiers principes}, 1.ª partie, ch. IV). Muito bem dito. --- D' onde se vê o erro de \emph{Locke} e dos \emph{sensualistas}, quando, confundindo o \emph{infinito}\pgn{342}com o \emph{indefinito}, disseram que o \emph{infinito} se forma pela somma dos \emph{finitos}. O \emph{infinito} exclue todo o limite. Ora o finito sommado com o finito dá sempre o finito, e não o infinito. Alem d'isso, se o \emph{infinito} fosse o resultado da somma dos entes finitos, seguir-se-hia que, subtrahindo, uma a uma, as partes sommadas, deveriamos chegar a exgottar o infinito, que por isso deixaria de ser infinito, pois este é necessariamente inexgottavel.\par Os principaes erros, relativos á \emph{origem} da idea do \emph{infinito}, podem reduzir-se a tres : ao dos \emph{Cartesianos} e dos \emph{innatistas}, ao dos sequazes do \emph{ontologismo}, e ao dos \emph{pragmatistas}. --- Os \emph{Cartesianos} e os \emph{innatistas} sustentam que a idea do infinito é infundida por \textsc{Deus} na intelligencia humana, e que ella é o fundamento de todas as nossas percepções. Mas não é assim; porque, como veremos, não se dão ideas \emph{innatas}, e, para julgarmos que um ente é \emph{limitado}, não é necessario recorrer ao ente \emph{illimitado}, mas basta considerar as creaturas, cada uma das quaes, por não ter as perfeições das outras, concebe-se como limitada. --- Os \emph{ontologistas} ensinam que percebemos o \emph{infinito} em si mesmo, por intuição, e que da idea do infinito deduzimos a do finito. Tambem esta opinião é falsa; pois ninguem tem consciencia d' essa intuição, ou, melhor, todos temos consciencia de que não percebemos o \emph{infinito} por intuição, porque, se assim o percebessemos, a idea que temos d' elle seria perfeitissima, clarissima : o que não é. --- Os \emph{pragmatistas} sustentam com \emph{Hügel} que o homem tem uma experiencia intima, verdadeira e real, embora indeterminada, do \emph{infinito}, e que por isso pelo sentimento religioso e pela experiencia religiosa percebemos, ainda que de um modo confuso, o \emph{infinito}. Tambem esta opinião é inacceitavel. É certo que a nossa alma, creada por Deus e por Elle elevada á ordem sobrenatural, tem uma aspiração irresistivel para o infinito, um desejo ardente de conhecer perfeitamente a Deus, mas tambem é certo que o sentimento religioso e a experiencia religiosa, sendo coisas sujeitas a enganos, não são por si sufficientes para nos dar uma idéa exacta do infinito, e por isso precisam do auxilio do superior criterio da intelligencia, que as regule. Alem d'isso, só uma faculdade dotada de força abstractiva pode formar a idea do \emph{infinito}, excluindo da entidade finita todo o limite; ora o sentimento religioso e a consciencia religiosa não têm essa força abstractiva indispensavel. Finalmente, se o sentimento religioso e a consciencia religiosa podessem attingir ou perceber directamente o \emph{infinito}, a idea \emph{do infinito} seria exclusivamente \emph{positiva}; ora não é assim (\emph{Cf}. \emph{Enc}. « \emph{Pascendi} » do S. P. Pio X, 8 de Set. de 1907).}
\pg{342} \textbf{65. O finito é cognoscivel em si mesmo.} --- O \emph{finito}, emquanto tal, importa duas coisas: a \emph{entidade} e a \emph{limitação}. Ora nós podemos muito bem conhecer, em si mesma, a \emph{entidade}, que é uma coisa absoluta, como podemos conhecer, na sua entidade, a \emph{limitação}, que é apenas a negação de ulterior perfeição. Logo o \emph{finito} é cognoscivel em si mesmo \footnote{\emph{Descartes}, como vimos, sustentou que a nossa intelligencia concebe primeiramente a idea do \emph{infinito}, e depois a do \emph{finito}; porque, na opinião d'elle, não poderiamos dizer que um ente é limitado, se não tivessemos já a idea do ente \emph{illimitado}. --- Mas é manifesto o seu erro. O \emph{finito} é a \emph{entidade limitada}. Ora, como dizemos no texto, tanto a \emph{entidade}, quanto a\pgn{343}\emph{limitação}, podem ser percebidas em si mesmas, sem que primeiramente percebamos o \emph{infinito}; porque a \emph{entidade} e a \emph{limitação} nem importam infinidade, nem se referem ao infinito.}.
\pg{343} \textbf{66 O ente, que é acto puro, é infinito; o ente, que é infinito, é acto puro.}

a) \emph{O ente, que é acto puro, é infinito.} --- O ente, que é \emph{acto puro}, tem o ser, não recebido na essencia, mas por si subsistente. Ora o ente, que tem o ser, não recebido na essencia, mas por si subsistente, é \emph{infinito}; pois a limitação do ser deriva da essencia, em que é recebido. Logo o ente, que é \emph{acto puro}, é \emph{infinito}.

b) \emph{O ente, que é infinito, é acto puro.} --- O ente \emph{infinito} é dotado de toda a perfeição possivel, é a propria perfeição. Ora o ente, dotado de toda a perfeição possivel, é necessariamente \emph{acto puro}; porque só se lhe poderia attribuir alguma potencialidade, se lhe faltasse alguma perfeição, mas então deixaria de ser infinito. Logo o ente, que é \emph{infinito}, é \emph{acto puro} (¹).

\textbf{67. O ente, que consta de acto e de potencia, é finito; o ente, que é finito, consta de acto e de potencia.}

a) \emph{O ente, que consta de acto e de potencia, é finito.} --- O ente, que consta de \emph{acto} e de \emph{potencia}, tem o proprio ser recebido na essencia. Ora o ser recebido na essencia é limitado pela propria essencia, em que é recebido. Logo o ente, que consta de \emph{acto} e de \emph{potencia}, \emph{é finito} \footnote{A verdade d'estas duas proposições apparece evidente pela simples declaração dos termos. \emph{Acto} é perfeição; acto \emph{puro} é perfeição pura, sem sombra de imperfeição, e por isso illimitada, infinita. Vice-versa, o ente \emph{infinito}, que é o proprio ser, é \emph{acto puro}; porque, se fosse composto de potencia e de acto, seria limitado, deixaria de ser infinito. --- Possuindo todas as perfeições possiveis, o infinito possue todas as perfeições dos entes finitos, tanto as \emph{simples} como as \emph{mistas}, --- aquellas de um modo \emph{formal} e \emph{eminente}, estas de um modo \emph{virtual} e \emph{eminente}. Por isso, os termos, que exprimem as perfeições \emph{simples}, quando attribuidos a Deus, não egualam nem limitam a coisa significada, como quando se attribuem ás creaturas, mas deixam-na incomprehensivel e superior á significação dos proprios termos. (\emph{Sum}. \emph{Th}., p. I, q. 13, a. 5).}.\footnote{Só o \emph{acto puro}, sendo o proprio ser subsistente, é \emph{infinito}. Onde se encontra a potencialidade, ahi se encontra a limitação. --- Alem d'isso, o \emph{infinito} não é o resultado da addição de partes. Portanto, onde se encontra a parte, ahi se encontra a limitação. Nem vale recorrer a um numero immenso de partes; o finito, addicionado ao finito, dá sempre o finito.\par D'ahi os seguintes corollarios da mais alta importancia\par 1.º) \emph{Não pode haver um infinito maior do que outro infinito.} --- O infinito não admitte o mais e o menos; e por isso, se fosse \emph{excedido} por outro infinito\pgn{344}, deixaria de ser tal. E não só não admitte um outro infinito maior, mas nem comporta outro \emph{egual}; porque, se fosse duplice ou multiplice, deixaria de ser infinito. Por isso, deve ser unico. --- Quando se diz que não pode haver um infinito maior que outro infinito, falla-se no infinito, que é tal \emph{em acto} ou \emph{absolutamente}, e não se falla no infinito, que só é tal \emph{em potencia}, emquanto é sempre capaz de receber novas partes, ou novas unidades, e que \emph{em acto}, ou \emph{na realidade}, é sempre finito. --- O infinito \emph{em potencia} não é infinito sob todos os aspectos nem em todas as ordens, mas é tal apenas sob um certo aspecto e numa só ordem; nem é uma substancia, mas é um attributo accidental das coisas, que se dizem infinitas, e por isso pode multiplicar-se, conforme se multiplica o sujeito, a que se attribue. Assim, podemos imaginar muitas linhas infinitas no comprimento, infinitas series de numeros pares e impares, infinitas series de effeitos possiveis em cada uma das causas, infinitos individuos possiveis em cada especie, etc. --- Como tambem, este infinito \emph{em potencia} não comporta outro infinito maior na mesma ordem, em que é infinito; mas, por causa de um sujeito maior, pode admittir outro infinito, maior do que elle, numa \emph{outra ordem}. Assim, as especies dos numeros pares são infinitas, e tambem infinitas são as especies dos numeros impares; e, todavia, a collecção dos numeros pares e impares é maior do que a dos numeros pares. Esta theoria é de S. Thomaz (\emph{Quodlib}. IX, q. 1, a. 1 ad 1).\par 2.º) \emph{Nenhuma creatura pode ser infinita quanto á essencia.} --- O ente, que é infinito quanto á \emph{essencia}, é o proprio Ser subsistente, acto purissimo, perfeição sem limite. Ora nenhuma creatura, sendo contingente, é o proprio ser subsistente, o acto purissimo, a propria perfeição. Logo nenhuma creatura pode ser infinita \emph{quanto á essencia}. Em menos palavras: creatura infinita \emph{na essencia} é uma contradicção. Por isso S. Thomaz nega que Deus possa produzir uma creatura infinita quanto á essencia. (\emph{Sum}. \emph{Th}., p. I, q. 4, a. 2 ad 1). --- Para a refutação das objecções, basta considerar --- a) que um effeito infinito é um effeito necessario, independente de todo e qualquer ente, e por isso é um effeito que não é effeito; --- b) que o effeito da causa infinita deve demonstrar e demonstra a infinidade da causa, --- não pela infinidade da propria essencia (o que é absurdo), --- mas pela infinidade da força necessaria para a sua producção, porque só uma força infinita pode crear, ou tirar do nada as coisas.\par 3.º) \emph{Nenhuma creatura pode ser a mais perfeita de todas.} --- \emph{Durando}, \emph{Henrique de Gand}, \emph{Gabriel Bie}, Occam, \emph{Arriaga}, e muitos \emph{Scotistas e Nominalistas} ensinaram que pode haver uma creatura tão perfeita, que não possa haver outra mais perfeita. \emph{S. Thomaz} e, geralmente, os \emph{Scholasticos} são contrarios a essa doutrina, e com razão. Na verdade, uma creatura, porquanto se supponha perfeita, nunca pode exgottar a participabilidade da perfeição divina, e entre essa creatura e Deus existe sempre uma distancia infinita, de modo que podem dar-se ainda infinitas outras creaturas, cada uma das quaes seja mais perfeita do que a que se diz perfeitissima, sem que nenhuma chegue a exgottar a divina participabilidade. Portanto, uma creatura, que seja a mais perfeita de todas, está fóra da serie dos possiveis. ---\pgn{345}Egualmente repugna uma creatura que seja a mais imperfeita de todas, de modo que não se conceba outra mais imperfeita. Toda a creatura, embora constituida numa especie imperfeitissima, comtudo dista infinitamente do nada, e por isso, abaixo d'ella, pode haver outras creaturas sem fim.\par 4.º) \emph{Nenhuma creatura pode ser dotada de força} ou \emph{actividade infinita.} --- É um corollario da primeira proposição. A força ou a actividade deriva da essencia, e é proporcionada a ella. Se a essencia das creaturas é necessariamente finita, tambem a sua força ou actividade é finita. (Cf. \emph{C. Gent}., I, 43). --- Por isso, a creatura racional, embora tenda, como para o seu fim, para o Bem infinito, comtudo só o alcança por um modo finito, isto é, por actos finitos de faculdades finitas, que são a intelligencia e a vontade.\par 5.º) \emph{Uma grandeza} ou \emph{extensão actualmente infinita repugna.} --- Uma essencia determinada é dotada de propriedades determinadas. Ora a grandeza ou extensão é uma propriedade do corpo, e todo o corpo tem uma essencia determinada. Logo a grandeza ou extensão do corpo é determinada, e por isso finita. D'onde se segue que a extensão de um corpo tem certos limites, que não é possivel transpor, sem que o proprio corpo seja physicamente destruido, embora a extensão, considerada em si, seja ulteriormente divisivel. --- O que dizemos da \emph{quantidade} pode applicar-se á \emph{qualidade}. Repugna que uma qualidade seja infinita na sua intensidade; pois, como dissemos, uma propriedade deve ser e é sempre proporcionada ao sujeito, de que deriva e em que reside. --- Todavia, a grandeza, considerada de um modo \emph{abstracto}, diz-se \emph{potencialmente} infinita, emquanto pode augmentar sem fim.\par 6.º) \emph{Uma multidão actualmente infinita repugna.} --- Toda a \emph{multidão} é uma especie de \emph{numero}, resulta da addição de unidades, e é sempre capaz de augmento e de diminuição. Ora o \emph{infinito} --- não pode ser uma especie de numero (pois o numero é uma multidão medida pela unidade, emquanto o infinito, por ser tal, não tem limites e por isso não pode ser medido), --- não resulta da addição de unidades (porque as unidades são partes finitas, e estas não podem dar o infinito), --- não é capaz de augmento e de diminuição (porque uma multidão, á qual pode accrescentar-se uma unidade, já se presuppõe finita; assim como uma multidão, que se suppõe infinita, se lhe se podesse tirar uma unidade, deixaria de ser infinita). Logo uma multidão actualmente infinita repugna.\par Dissemos --- \emph{actualmente} infinita, para se denotar que a multidão infinita repugna, quando todos os seus elementos se consideram \emph{em acto}, isto é, na ordem \emph{da existencia real}; porque, se os elementos se consideram \emph{em potencia}, isto é, na ordem \emph{da possibilidade}, neste caso a multidão infinita não repugna. Uma multidão infinita de entes \emph{possiveis} é possivel. Na verdade, o augmento de uma multidão está na razão directa da divisão da grandeza ou extensão; visto que, quanto mais se divide uma grandeza, tanto mais cresce o numero das partes. Ora, assim como a grandeza é divisivel sem fim, assim tambem a multidão cresce sem fim, isto é, a infinidade \emph{potencial} da multidão é uma consequencia da infinidade \emph{potencial} da grandeza ou extensão. Logo uma multidão infinita \emph{em potencia} é possivel. (\emph{Sum}. \emph{Th}., q. 1, q. 7, a. 4). --- Agora uma advertencia. A multidão, que se diz infinita \emph{em potencia}, emquanto pode ser dividida em partes, cada vez mais pequenas, \emph{ou}\pgn{346}pode receber outras unidades, sem nunca chegar ao fim, é, de facto, sempre limitada, porque constituída por um numero finito de unidades. Por isso o infinito \emph{em potencia} costuma chamar-se, mais propriamente, \emph{indefinito}.}

\pg{344}

b) O \emph{ente}, que é \emph{finito}, consta \emph{de acto} e \emph{de potencia}. --- O ente, que é \emph{finito}, tem o proprio ser limitado. Ora o ser emtanto é \pg{345}limitado, emquanto é recebido numa determinada essencia, como\pg{346} \emph{acto} na \emph{potencia}. Logo o ente, que é \emph{finito}, consta de acto e de \emph{potencia} \footnote{O \emph{ser}, por si, significa \emph{acto}, \emph{perfeição}, e por isso exclue do seu conceito a \emph{limitação}, que importa imperfeição. Se, pois, o ser se acha limitado, a unica razão é porque é recebido, como acto na potencia, numa certa essencia, a qual tem uma capacidade finita. D'ahi a composição do ente finito.}.

\textbf{68. O ente infinito é necessario; o ente necessario é infinito.}

a) \emph{O ente infinito é necessario}. --- Ente \emph{infinito} é o que encerra todas as perfeições possiveis, e por isso é o proprio ser subsistente, cuja existencia é uma e a mesma coisa com a essencia. Ora um ente, cuja existencia é uma e a mesma coisa com a essencia, existe pela exigencia da sua essencia, não pode ser privado da existencia, sem ser privado da essencia, e por isso é \emph{necessario}. Logo o ente \emph{infinito é necessario}.

b) O \emph{ente necessario é infinito}. --- Ente \emph{necessario} é o que existe pela exigencia da sua natureza, é o proprio ser subsistente, acto purissimo, sem sombra de potencialidade. Ora um ente que é o proprio ser subsistente, acto purissimo, sem sombra de potencialidade, encerra todas as perfeições possiveis, e por isso é \emph{infinito}; pois; se não fosse infinito, por lhe faltar alguma perfeição, seria em potencia para ella e deixaria de ser acto purissimo. Logo o ente \emph{necessario} é \emph{infinito} ().\footnote{Tambem a verdade d'estas proposições se conhece pela explicação dos termos. --- \emph{O} ente \emph{infinito} possue todas as perfeições possiveis e por isso exclue toda e qualquer potencialidade, que é imperfeição. Excluindo toda e qualquer potencialidade, \emph{o} ente \emph{infinito} exclue tambem a potencialidade para a não-existencia, e por isso não pode deixar de existir, é absolutamente \emph{necessario}. --- Como tambem o ente \emph{necessario} é existente por si mesmo, pela exigencia da sua essencia. Sendo existente por si mesmo, o ente \emph{necessario} é acto purissimo, e portanto encerra todas as perfeições possiveis, é \emph{infinito}. Com effeito, o ente necessario possue o ser pela exigencia da propria essencia. Ora o que um ser possue pela sua essencia não pode ser diminuto ou deficiente, porque a essencia não comporta nem o mais nem o menos, mas é \emph{indivisivel}, e por isso o que se possue pela exigencia da essencia, possue-se de um modo completo e illimitado. \textsc{Diz-S}. Thomaz: «Quod alicui convenit ex sua natura, et non ex aliqua causa, minoratum in eo et deficiens esse non potest. Si enim naturae aliquid essentiale subtrahitur vel additur, \emph{iam altera natura} erit» (\emph{C}. \emph{Gent}, II, 15). Logo o ente necessario possue o ser e a perfeição de um modo illimitado e por isso é a plenitude do ente e da perfeição. Quanta luz irradia dos principios do Angelico Doutor!}
\pg{347} \textbf{69. O ente finito é contingente; o ente contingente é finito.}

a) \emph{O ente finito é contingente}. --- Ente \emph{finito} é aquelle que tem o ser limitado. Ora o ente, que tem o ser limitado, não é o proprio ser subsistente, não existe pela exigencia da sua essencia, é \emph{contingente}. Logo o ente \emph{finito} é \emph{contingente}.

b) \emph{O ente contingente é finito}. --- Ente \emph{contingente} é o que é por si indifferente para a existencia. Ora o ente, que é por si indifferente para a existencia, é composto de potencia e de acto, e por isso é \emph{finito}. Logo o ente \emph{contingente} é \emph{finito} \footnote{\emph{Finito} e \emph{contingente} são termos, que se correspondem. --- O \emph{finito} é composto de acto e de potencia (pois a limitação é potencia, é imperfeição), não é o proprio ser subsistente, não existe pela exigencia da sua essencia, é indifferente para a existencia e para a não-existencia, é \emph{contingente}. --- O \emph{contingente}, pela sua vez, é composto de essencia e de existencia, de potencia e de acto, e por isso é \emph{finito}; pois cada parte é finita, e, como o finito, addicionado ao finito, só dá o finito, a somma das partes é tambem finita.}.

70. A existencia dos entes finitos suppõe a \textbf{existencia} do Ente infinito.

a) \emph{Os entes finitos}, porque são \emph{entes}, possuem o ser, e, porque são \emph{finitos}, não são o proprio ser, isto é, não possuem o ser por si mesmos, pela exigencia da sua essencia. Não possuindo o ser por si mesmos, os entes finitos devem tel-o recebido de outro ente; aliás haveria um effeito sem causa proporcionada. Esse outro ente, ou é tambem finito, ou é infinito. Se é infinito, temos o que queriamos. Se é finito, subiremos \textbf{ainda, até se encontrar} necessariamente o ente, que possua o \textbf{ser} \textbf{por si mesmo}, que seja a plenitude do ser, do qual todos os \textbf{entes finitos} receberam o ser, que seja \emph{infinito}; porque, se não existisse o ente infinito, não poderia existir nenhum ente finito. Logo a existencia dos entes \emph{finitos} suppõe a existencia do Ente \emph{infinito} \footnote{O argumento é claro e convincente. Todo o ente, que existe, --- ou existe \emph{por si}, pela sua essencia, e é \emph{ens a se}. --- ou existe \emph{por outro ente}, por participação, e é \emph{ens ab alio}. Não se dá meio termo. Ora o ente \emph{finito}, como tal, não é ente \emph{por si}, isto é, não existe pela exigencia da sua essencia, mas é ente por participação, porque a sua essencia não é a sua existencia, não é o seu ser. Logo existe por outro ente, que lhe communicou a existencia. Este outro ente deve ser, em ultima analyse, o proprio ser subsistente, o ente infinito. Se não existisse o infinito, o ente por essencia, não poderia existir o finito, o ente por participação: assim como, se não existisse o calor, não poderia existir nenhum corpo quente, e, se não existisse a luz, não poderia existir nenhum corpo luminoso. Diz S. Thomaz: «Quod per essentiam dicitur est causa omnium, quae per participationem dicuntur, sicut\pgn{348}ignis est causa omnium ignitorum in quantum hujusmodi. Deus autem est ens per essentiam suam, quia est ipsum esse; omne autem aliud ens est ens per participationem, quia ens quod sit suum esse non potest esse nisi unum. Deus igitur est causa essendi omnibus aliis» (\emph{C}. \emph{Gent}., II, 15).}.
\pg{348} b) Os entes finitos, ainda que todos possuam uma realidade ou perfeição limitada, todavia nem todos a possuem no mesmo grau, mas um é mais perfeito, e outro o é menos; assim o ente vivo é mais perfeito do que o ente não-vivo. Tal desegualdade deve ter a sua causa sufficiente. Esta causa deve encontrar-se --- ou nos proprios entes finitos, --- ou num ente infinito. Ora não pode encontrar-se nos entes finitos; porque estes, sendo contingentes, não têm em si a razão da sua existencia, e por isso nem do grau da sua realidade ou perfeição. Logo deve encontrar-se num ente infinito, que, espalhando diversamente as perfeições em conformidade com os archetypos da sua intelligencia, produziu a maravilhosa variedade do universo. Logo a existencia dos entes \emph{finitos} suppõe a existencia do Ente \emph{infinito} \footnote{No mundo ha muitas e diversas perfeições, mas todas limitadas e num grau desigual; assim, no mesmo genero de seres vivos, o homem é mais perfeito do que o animal, e este é mais perfeito do que a planta. Ora tal limitação e desigualdade, que se encontra na perfeição das creaturas, não pode derivar da essencia ou natureza d'ellas. Na verdade, segundo S. Thomaz, uma perfeição, que convém a um ente pela sua essencia, não pode ser de modo algum limitada, ou diminuta, sem que a essencia deixe de ser o que é, mas deve existir nelle no grau mais elevado, de que a propria perfeição é susceptivel. Logo essa limitação e desigualdade deve derivar, em ultima analyse, de uma causa extrinseca, que possua por si todas as perfeições num grau eminente e que as communique por differentes modos aos entes creados (\emph{Gent.}, II, 15). --- O S. Doutor ensina tambem que, quando uma coisa se diz e é mais perfeita que outra, essas coisas dizem-se e são mais ou menos perfeitas pelo facto de dependerem e de se approximarem, mais ou menos, de um ente unico e perfeitissimo: porque, se cada coisa possuisse \emph{por si} mesma a perfeição, de que é dotada, não haveria razão por que uma devesse ser mais e outra menos perfeita (\emph{De pot.}, q. 3, a. 5).}.

71. \textbf{Corollarios.} --- Das doutrinas expostas resulta :

I. O infinito presuppõe-se ao finito --- como \emph{causa}, porque o infinito produz o finito, --- mas não como \emph{objecto}, porque o infinito é por nós conhecido depois e por meio do finito\footnote{A existencia \emph{real} do ente \emph{infinito} não pode ser demonstrada apenas pelo \emph{conceito} do mesmo ente. É certo que nada pode conceber-se que seja mais perfeito do que o ente \emph{infinito}, e que a existencia é uma perfeição; mas, emquanto não se sahe do conceito, a existencia, devida ao Ente infinito, será apenas \emph{ideal}. É o que dissemos em relação ao ente \emph{necessario}.}.
\pg{349} II. Dada a hypothese da existencia real de um só ente finito, deduz-se legitimamente a existencia do Ente infinito \footnote{Ente \emph{simples} é o que não tem partes. --- Se se considera a razão ou o fundamento da \emph{simplicidade} do ente, este pode ser, como dizemos no texto, \emph{negativo}, \emph{positivo}, e \emph{abstractivo} ou \emph{precisivo}. --- Se se considera a qualidade das partes, que se excluem do ente \emph{simples}, este divide-se principalmente, e em opposição ao \emph{composto}, em \emph{physico}, \emph{metaphysico} e \emph{logico}, conforme as partes, que se excluem, são \emph{physicas} (materia e forma), ou \emph{metaphysicas} (potencia e acto, essencia e existencia), ou \emph{logicas} (genero e differença). --- Se se considera o grau da \emph{simplicidade} do ente, este, como dizemos no texto, pode ser \emph{absoluto} e \emph{relativo}, conforme exclue --- toda e qualquer especie de \emph{composição} (\textsc{Deus}, é só Elle, é \emph{absolutamente} simples), --- ou somente uma ou outra especie (a nossa alma é \emph{relativamente} simples, porque, se exclue a composição \emph{physica}, não exclue a \emph{metaphysica}).}

ARTIGO III

\emph{Ente simples e ente composto}.

72. Ente simples e ente composto. --- Ente \emph{simples} é o que não tem partes, e por isso é indivisivel. Pode ser \emph{negativo}, \emph{positivo} e \emph{abstractivo}, conforme pode não ter partes --- ou pela tenuidade e imperfeição da sua entidade, como é o ponto mathematico, --- ou pela perfeição do seu ser, como é a nossa alma, --- ou pelo effeito da abstracção, como é o ente em geral. Neste artigo occupamo-nos do \emph{simples positivo}. --- Ente \emph{composto} é o que tem partes, e por isso é divisivel. Pode ser \emph{physico}, \emph{metaphysico} e \emph{logico}, conforme as partes, são --- \emph{physicas}, como a alma, e o corpo, --- ou \emph{metaphysicas}, como a potencia e o acto, a essencia e a existencia, --- ou \emph{logicas}, como o genero e a differença.

73. \textbf{O ente simples, quando subsiste perfeitamente, é mais perfeito que o composto.} --- O ente \emph{simples}, que subsiste perfeitamente, por não constar de partes, tem todos os elementos ou graus, que constituem a sua entidade, não só unidos, mas indivisiveis; ao passo que o ente \emph{composto}, embora tenha os seus elementos unidos, comtudo não os tem indivisiveis. Ora um ente,

(1) Nem se diga que entre o finito e o infinito não existe proporção, e por isso o conhecimento d'aquelle não pode levar ao conhecimento d'este. Por quanto, se entre o finito e o infinito não existe uma proporção de \emph{entidade} e de \emph{egualdade}, existe uma proporção de \emph{dependencia} e de \emph{analogia}; pois o finito depende do infinito, como o effeito depende da causa. --- Esta falta de proporção de \emph{entidade} e de \emph{egualdade} obsta a que tenhamos um conhecimento proprio e adequado do infinito; mas não obsta a que conheçamos a existencia do mesmo infinito, e, de algum modo, as suas perfeições.
\pg{350} que tem os elementos da sua entidade unidos e indivisiveis, é mais perfeito que o ente que os tem unidos mas divisiveis; porque, ao passo que o primeiro não pode perder a sua unidade e, com a unidade, a entidade, o segundo pode perder uma e outra coisa (o que é imperfeição) Logo o ente \emph{simples}, quando subsiste perfeitamente, é mais perfeito que o \emph{composto}.

Dizemos --- quando \emph{subsiste perfeitamente}; porque, se o ente \emph{simples}, ou não \emph{subsiste}, ou não subsiste \emph{perfeitamente}, mas tem razão de \emph{parte}, é menos perfeito que o \emph{composto}. Assim a alma do animal, embora simples, é comtudo menos perfeita que o animal inteiro; e a nossa alma, embora seja simples e subsistente, comtudo, por não subsistir \emph{perfeitamente} quanto á \emph{especie}, é menos perfeita que o composto humano \footnote{A \emph{simplicidade}, considerada em si, emquanto é negação de composição, não denota necessariamente nem perfeição nem imperfeição: porque, diz S. Thomaz, o que constitue a excellência do ente não é propriamente a simplicidade, mas é a perfeição, é a somma dos graus de entidade, real ou logicamente distinctos, de que o ente é constituido. Para se averiguar, pois, em concreto, se a \emph{simplicidade} importa perfeição ou imperfeição, é necessario considerar a natureza do ente, e a razão e o modo da sua simplicidade. Ora o ente pode dizer-se e ser \emph{simples}, --- ou emquanto importa um \emph{todo completo} e \emph{subsistente} na mesma simplicidade, --- ou emquanto importa apenas uma \emph{parte}, que não se divide noutras partes. No primeiro caso, a \emph{simplicidade} denota \emph{perfeição}, e o ente \emph{simples}, que é um todo completo e subsistente, é mais perfeito que o \emph{composto}; porque, como dizemos no texto, conserva melhor a sua unidade e por isso a sua entidade. No segundo caso, a \emph{simplicidade} denota \emph{imperfeição} (porque quem diz \emph{parte} diz limitação) e o ente \emph{simples} é menos perfeito que o \emph{composto}. (²) O Angelico Doutor escreve: « \emph{Simplicitas} per se non est causa nobilitatis, sed \emph{perfectio}; unde ubi perfecta bonitas \emph{in} uno simplici invenitur, simplex est nobilius quam compositum: quando autem e contrario simplex est imperfectum, tunc compositum est nobilius quam simplex » (In IV \emph{Sent}., dist. 11, q. 2, a. 1 ad 1).\par Portanto a \emph{simplicidade} do ente subsistente não exprime apenas uma negação de composição, mas denota uma perfeição positiva: pois denota a entidade emquanto subsiste na simplicidade. Assim como a unidade não indica apenas a negação de divisão, mas significa o proprio ente, do qual exclue a divisão; assim tambem a simplicidade significa por si a propria perfeição positiva da entidade, da qual exclue a composição. Por isso diz S. Thomaz que, no mesmo genero, um ente é tanto mais excellente, quanto é mais simples: « In quolibet genere, tanto aliquid est \emph{nobilius}, quanto \emph{simplicius} » (\emph{C}. \emph{Gent}., I, 18). Ora a excellencia indica uma perfeição positiva.}.

\textbf{74. O ente, que é acto puro, é absolutamente simples; o ente, que é absolutamente simples, é acto puro.}

a) \emph{O ente, que é acto puro, é absolutamente simples.} --- O\pg{351} ente, que é \emph{acto puro}, não encerra em si mesmo nenhuma potencialidade. Ora o ente, que não encerra em si nenhuma potencialidade, é absolutamente \emph{simples}; porque é só acto, e o acto não tem partes, em que possa dividir-se. Logo o ente, que é \emph{acto puro}, é absolutamente \emph{simples}.

b) \emph{O ente, que é absolutamente simples, é acto puro.} --- O ente, que é absolutamente \emph{simples}, não consta de partes, e por isso não é composto de potencia e de acto. Ora o ente, que não é composto de potencia e de acto, é \emph{acto puro}. Logo o ente, que é absolutamente \emph{simples}, é \emph{acto puro} \footnote{O \emph{acto puro} é absolutamente \emph{simples}. O acto é perfeição, e o acto puro é só perfeição, é perfeição subsistente; e por isso exclue toda a potencialidade, não pode ser composto de potencia e de acto, é absolutamente simples. Por quanto, como diz S. Thomaz, em todo o composto deve encontrar-se o acto e a potencia. Com effeito, muitas coisas não podem tornar-se simplesmente uma só coisa, se não emquanto uma seja acto e outra potencia; visto que, se fossem todas em acto, haveria um aggregado de coisas, mas nunca uma só coisa composta. Ora o acto puro, excluindo toda e qualquer potencialidade, é absolutamente simples. --- Não é menos evidente que o ente absolutamente \emph{simples} é \emph{acto puro}. Por quanto, o ente absolutamente \emph{simples}, que subsiste perfeitamente em si mesmo, ou é só potencia, ou é só acto. Ora é absurdo dizer que é só potencia; pois, como vimos, a potencia, denotando limitação e imperfeição, não pode subsistir em si mesma. Logo deve ser só acto, deve ser acto puro (\emph{Cf}. \emph{C}. \emph{Gent}., I, 18).}.

\textbf{75. O ente composto não pode ser acto puro.} --- Ente \emph{composto} é o que tem partes. Ora o que tem partes não pode ser \emph{acto puro}; porque --- ou uma das partes será \emph{acto} em relação a outra, --- ou todas as partes serão como que \emph{em potencia} em relação ao todo. Logo o ente \emph{composto} não pode ser \emph{acto puro} \footnote{Diz S. Thomaz: «In omni composito oportet esse potentiam et actum; quia vel una pars est actus respectu alterius, vel saltem omnes partes sunt in potentia respectu totius» (\emph{Sum. Th.}, p. I, q. 3, a. 7).}.

\textbf{76. O ente, que é absolutamente simples, é necessario e infinito; o ente, que é necessario e infinito, é absolutamente simples.}

a) \emph{O ente, que é absolutamente simples, é necessario e infinito.} --- O ente, que é absolutamente \emph{simples}, é acto puro, é o proprio ser subsistente. Ora o acto puro, o proprio ser subsistente, é \emph{necessario} e \emph{infinito}. Logo o ente que é absolutamente \emph{simples}, é \emph{necessario} e \emph{infinito}.

b) \emph{O ente, que é necessario e infinito, é absolutamente simples.} --- O ente, que é \emph{necessario} e \emph{infinito}, possue a existencia\pg{352} identificada com a essencia e é acto puro, sem sombra de potencialidade. Ora o ente, que possue a existencia identificada com a essencia e que é acto puro, não é composto de potencia e de acto, mas é absolutamente \emph{simples}. Logo o ente, que é \emph{necessario} e \emph{infinito}, é absolutamente \emph{simples} \footnote{A verdade d'estas proposições é tambem manifesta pela simples inspecção dos termos. S. Thomaz prova que o ente \emph{necessario} é \emph{simples}, porque o ente \emph{composto}, emquanto \emph{composto}, pode estar sujeito á dissolução, e por isso pode deixar de existir. (\emph{C}. \emph{Gent}., I, 18).}.

77. O ente, que é composto, é contingente e finito; o ente, que é contingente e finito, é composto.

a) \emph{O ente, que é composto, é contingente e finito.} --- O ente, que é \emph{composto}, não é puro acto, mas resulta de potencia e de acto. Ora o ente, que resulta de potencia e de acto, --- é \emph{contingente}, porque não existe pela exigencia da sua essencia, mas é indifferente quanto á existencia, pode ser e deixar de ser, --- e é \emph{finito}, porque é composto de partes, essencialmente finitas. Logo o ente, que é \emph{composto}, \emph{é contingente e finito}.

b) \emph{O ente, que é contingente e finito, é composto.} --- O ente, que é contingente e finito, possue a existencia realmente distincta da essencia e resulta de varias partes. Ora o ente, que possue a existencia realmente distincta da essencia e resulta de varias partes, é \emph{composto}; porque consta de potencia e de acto. Logo o ente, que é \emph{contingente} e \emph{finito}, \emph{é composto} \footnote{O \emph{composto} é \emph{contingente}; porque, como diz S. Thomaz, o \emph{composto}, emquanto tal, está sujeito á dissolução e por isso pode deixar de ser. --- Noutro logar o S. Doutor escreve: « Omne id, in cujus substantia admiscetur potentia, secundum id quod habet de potentia, potest non esse; quia quod potest esse, potest non esse » (\emph{C}. \emph{Gent}., I, 16).}.

78. \textbf{O ente composto depende de uma causa extrinseca.} --- O ente \emph{composto} resulta da união de varias partes. A união, por ser um effeito, exige uma causa. Esta causa não pode ser o proprio \emph{composto}; porque, aliás, este, sendo naturalmente posterior ás partes componentes, operaria antes de existir e seria a causa de si mesmo. Deve ser, pois, uma causa \emph{extrinseca}. Logo o ente \emph{composto} depende de uma causa \emph{extrinseca} \footnote{O ente \emph{composto} é naturalmente posterior ás partes, de que resulta e que são diversas entre si, e por isso é um effeito. Qual é a causa d'este effeito? --- São as proprias partes? Não; porque \emph{partes diversas}, --- emquanto \emph{partes}, estão todas em potencia com relação ao todo, --- e, emquanto \emph{diversas}, não podem por si formar a união. --- É o proprio composto? Tambem não; \pgn{353}porque, então, o composto seria causa de si mesmo, e existiria antes de existir. Logo deve ser uma causa \emph{extrinseca}. Diz S. Thomaz: « Omne compositum causam habet. Quae enim secundum se diversa sunt, non conveniunt in aliquod unum, nisi per aliquam causam adunantem ipsa » (\emph{Sum}. \emph{Th}., P. I, q. 3, a. 7).}.
\pg{353} 79. O ente composto depende, immediata ou mediatamente, do ente simples. O ente \emph{composto} depende, como dissemos, de uma causa extrinseca. Esta causa --- ou é um ente \emph{simples}, --- ou é um ente \emph{composto}. Se é um ente \emph{simples}, está provada a conclusão. Se é um ente \emph{composto}, renovamos a pergunta, até encontrarmos o ente \emph{simples}; porque, se não existisse o ente \emph{simples}, não poderia existir nenhum ente \emph{composto}. Logo o ente \emph{composto} depende, immediata ou mediatamente, do ente \emph{simples} \footnote{O ente \emph{composto} depende do \emph{simples}, enquanto se resolve, em ultima analyse, em elementos \emph{simples}; e, tambem e sobretudo, emquanto exige uma causa distincta, que una as partes componentes e que deve encontrar-se finalmente no Ente subsistente na sua simplicidade. S. Thomaz prova esta proposição, como tambem a antecedente, com as seguintes palavras: « Omnis compositio indiget aliquo componente. Si enim compositio est, ex pluribus est. Quae autem secundum se sunt plura, in unum non conveniunt, nisi ab aliquo componente uniantur. Si igitur compositus esset Deus, haberet componentem; non enim ipse seipsum componere posset, quia nihil est causa sui ipsius, esset enim prius seipso, quod est impossibile» (\emph{C. Gent.}, I, 18). E mais abaixo diz: « Ante multitudinem oportet invenire unitatem. In omni autem composito est multitudo. Igitur oportet id, quod est ante omnia, scilicet Deum, omni compositione carere » (\emph{Ibid}.).}.

\textbf{80. Corollarios.} --- Dos principios expostos resulta:

I. Na ordem \emph{ontologica}, ou da \emph{realidade}, o \emph{simples} precede o \emph{composto}, porque este depende d'aquelle; mas, na ordem \emph{logica}, ou do \emph{conhecimento}, o \emph{composto} precede o \emph{simples}, porque definimos o \emph{simples} pelo \emph{composto}, que se apresenta primeiro á nossa intelligencia, e dizemos que o \emph{simples} é o ente destituido de partes, é o \emph{não-composto} (\textasciicircum{}).

\footnote{Devemos fazer outras considerações, relativas ao conhecimento do \emph{simples} e do \emph{composto}. --- O \emph{composto} pode ser conhecido quanto a uma ou mais partes e desconhecido quanto a outras: d'ahi os conhecimentos imperfeitos, relativos, constituidos por verdades e erros. --- \emph{O simples}, sendo uma coisa indivisivel, ou é conhecido totalmente, ou é totalmente desconhecido. Todavia, como o que é \emph{simples} em si mesmo, pode ser por nós considerado como uma coisa \emph{composta} e cognoscivel por varios juizos, alguns d'esses juizos podem ser verdadeiros e outros podem ser falsos. Assim Deus pode ser conhecido como Providencia e desconhecido como Espirito purissimo, e vice-versa. Se, porém, o \emph{simples} é tambem \emph{tal} para a nossa intelligencia, como são tantas verdades \emph{geometricas}, \emph{etc.}, então \emph{elle} é ou inteiramente conhecido ou inteiramente \emph{ignorado}, \emph{sem} a \emph{possibilidade de} um termo medio.\par}
\pg{354} II. Dada a existencia de um só ente \emph{composto}, deve admitir-se a existencia de um Ente absolutamente \emph{simples} (¹).

\textsc{Artigo} IV

\emph{Ente immutavel e ente mutavel}.

\textbf{§1. Ente immutavel é ente mutavel.} --- Ente \emph{immutavel} é o que não está sujeito á mudança, isto é, que não pode passar de um estado para outro, da potencia para o acto. --- Ente \emph{mutavel} é o que está sujeito á mudança, isto é, que pode passar de um estado para outro, da potencia para o acto\footnote{O homem, sendo um ente \emph{composto}, conhece primeiramente os entes \emph{compostos}, e por estes os \emph{simples}. Por isso, se não existisse nenhum ente \emph{composto}, o homem não poderia conhecer nenhum ente \emph{simples}, que não é objecto. proporcionado para nós. Mas, se existir um só ente \emph{composto}, podemos e devemos concluir que existe o Ente absolutamente \emph{simples}; porque o \emph{composto}, como dissemos, não pode ser o ente primeiro, mas suppõe e exige uma causa extrinseca, que, em ultima analyse, deve ser simples.}.

\textbf{§2. Especies da immutabilidade.} --- A \emph{immutabilidade} pode ser \emph{absoluta} e \emph{relativa}. --- É \emph{absoluta}, quando o ente exclue toda e qualquer mudança, quer esta derive de um principio intrinseco, quer derive de um principio extrinseco; tal é a \emph{immutabilidade} de \textsc{Deus}. --- É \emph{relativa}, quando o ente só exclue uma ou outra especie de mudança; tal é a \emph{immutabilidade} das creaturas, que, se são immutaveis quanto á \emph{essencia}, não são immutaveis quanto á \emph{existencia}\footnote{A \emph{immutabilidade} é uma consequencia da simplicidade, como a \emph{mutabilidade} o é da composição. A mudança importa uma passagem da potencia para o acto; por isso, onde se encontra a potencia, ahi se encontra a possibilidade da mudança, --- e, onde não se encontra a potencia, ahi se encontra a immutabilidade. --- A \emph{mudança} é, como dissemos, a passagem de um estado para outro. A \emph{mutabilidade} é a disposição do ente para passar de um estado para outro. Diz S. Thomaz: «De ratione \emph{mutationis} est quod aliquid idem se habeat \emph{aliter nunc} et \emph{prius}» (\emph{Sum}. \emph{Th}., p. I, q. 45, a. 2 ad 2).}.

\textbf{§3. Elementos da mudança.} --- Toda a mudança, por ser uma passagem da potencia para o acto, de um opposto para\footnote{Toda a mudança deriva --- ou de um principio intrinseco ao sujeito, que se muda, --- ou de um principio extrinseco, de que o sujeito mutavel depende. Se, portanto, um ente é \emph{immutavel} emquanto exclue toda e qualquer potencialidade, tanto em relação a si mesmo, quanto em ordem ás causas extrinsecas, a sua immutabilidade é \emph{absoluta}. --- Se o sujeito, considerado sob um certo aspecto, exclue a possibilidade da mudança, mas, considerado sob um outro aspecto, encerra essa possibilidade da mudança, a sua immutabilidade é \emph{relativa}.}\pg{355} outro, exige três elementos: --- 1º.) um \emph{principio}, de que se parte e que importa alguma coisa, a qual se perde; --- 2º.) um \emph{termo}, a que se chega e que importa alguma coisa, a qual se adquire; --- 3º.) um \emph{sujeito}, que do \emph{principio} chega ao \emph{termo} e se encontra num e noutro extremo. --- Se faltar algum d' estes elementos, não haverá verdadeira \emph{mudança} \footnote{O \emph{principio}, de que se parte, é chamado, pelos escholasticos, termo a \emph{quo}; o \emph{termo}, a que se chega, termo \emph{ad} quem. --- Entre o \emph{principio} e o \emph{termo}, ou entre o termo a \emph{quo} e o termo \emph{ad} quem, deve existir uma certa opposição; aliás não poderia conceber-se a mudança. Diz S. Thomaz: « \emph{Omnis mutatio est inter opposita} » (\emph{In V Physic.}, l. 2). --- A mudança recebe a espécie e a denominação --- não do \emph{principio}, ou do termo a \emph{quo}, --- mas do \emph{termo}, ou do termo \emph{ad quem}, porque é para este termo que a mudança tende e é nelle que se completa; assim dizemos que a água está quente, quando de fria se muda para quente. Escreve o S. Doutor: « \emph{Mutationes accipiunt speciem et dignitatem non a termino a quo, sed a termino ad quem} » (\emph{Sum. Th.}, p. I, q. 45, a. 1 ad 2). --- O \emph{sujeito} da mudança deve ser um e o mesmo no \emph{principio} e no \emph{termo}; o que muda propriamente é o \emph{estado} do sujeito. Esta mudança do estado realisa-se emquanto o sujeito perde uma coisa e adquire outra. --- Se, pois, o sujeito não se encontrasse no \emph{principio}, ou no \emph{termo}, não haveria propriamente \emph{mudança}. Assim não é mudança --- nem a \emph{creação}, porque antes d' ella não existe o sujeito, isto é, não há um sujeito commum aos dois extremos, --- nem a \emph{annihilação}, porque depois d' ella não existe o sujeito, isto é, não há um sujeito commum aos dois extremos, pois o \emph{nada} não é um \emph{estado}, mas uma mera negação (\emph{Cf. C. Gent.}, II, 17; \emph{De Pot.}, q. 3, a. 2).}.

\textbf{§4. Espécies de mudança.} --- A \emph{mudança} é --- \emph{intrinseca} ou \emph{extrinseca}, --- \emph{instantanea} ou \emph{successiva}, --- \emph{substancial} ou \emph{accidental}.

a) \emph{Mudança intrinseca} ou \emph{extrinseca}. --- É \emph{intrinseca}, quando o sujeito adquire ou perde alguma coisa em si mesmo; tal é a mudança de um corpo, que de frio se torna quente. --- É \emph{extrinseca}, quando o sujeito nada adquire ou perde em si mesmo, mas apenas reveste uma nova relação por se ter mudado um outro ente, que de algum modo lhe se refere; assim dizemos que uma torre, que estava á nossa direita, se muda emquanto passa para a nossa esquerda pelo facto de termos nós mudado de posição \footnote{Esta distinção não deve confundir-se com a distinção da mudança proveniente de uma causa \emph{intrinseca}, ou de uma causa \emph{extrinseca}. --- Mudança proveniente de uma causa \emph{intrinseca} é a que se funda na própria natureza de sujeito, que tem a potencialidade para a não-existencia, como é toda a substancia composta de materia e de forma. --- Mudança proveniente de uma causa \emph{extrinseca} é a que tem por \emph{principio} a potencia activa de um agente exterior, de que o sujeito depende; assim a nossa alma, embora não tenha\pgn{356}em si a potencia para a não-existencia, comtudo, por estar sujeita ao poder omnipotente de Deus, pode mudar. --- Como se vê, a mudança \emph{intrinseca} pode derivar de uma causa \emph{extrinseca}.}
\pg{356} b) \emph{Mudança instantanea} ou \emph{successiva}. --- São duas especies da mudança \emph{intrinseca}. --- É \emph{instantanea}, quando se effectua \emph{num instante}, e verifica-se quando o sujeito passa de um para outro estado \emph{contradictorio}, pois as coisas \emph{contradictorias} não admittem meios termos; tal é a mudança de \emph{não-homem} para \emph{homem}. --- É \emph{successiva}, quando se realisa \emph{successivamente}, e dá-se quando o sujeito passa de um para outro estado \emph{contrario}, pois as coisas \emph{contrarias} admittem varios meios termos; tal é a mudança da agua de fria para quente \footnote{A mudança \emph{successiva} chama-se propriamente \emph{movimento}, emquanto a \emph{instantanea} conserva o nome de \emph{mudança}. --- Embora a \emph{mudança} e o \emph{movimento} se empreguem indifferentemente, comtudo differem uma do outro. Com effeito, --- ao passo que o \emph{movimento} exige extremos \emph{positivamente} contrarios, a \emph{mudança} pode effectuar-se entre um extremo negativo e outro positivo, p. ex., entre a \emph{privação} e a \emph{forma}; --- ao passo que o sujeito do \emph{movimento} é um ente dotado de ser substancial, o sujeito da \emph{mudança} pode ser um ente privado do ser substancial, como é a \emph{materia prima}.}.

c) \emph{Mudança substancial} ou \emph{accidental}. --- São duas outras especies da mudança \emph{intrinseca}. --- É \emph{substancial}, quando o sujeito passa de uma para outra \emph{natureza}, ou \emph{essencia}; tal é a mudança do alimento, que se torna carne. --- É \emph{accidental}, quando o sujeito, conservando a sua natureza, passa de uma para outra \emph{propriedade accidental}; tal é a mudança do homem, que de ignorante se torna sabio.

\textbf{§5. O ente, que é acto puro, é immutavel; o ente, que é immutavel, é acto puro.}\footnote{A mudança \emph{substancial} divide-se em \emph{geração} e \emph{corrupção}. --- \emph{A geração} é a passagem do não-ser substancial para o ser substancial, como quando do não-homem se faz o homem. A \emph{corrupção} é a passagem do ser substancial para o não-ser substancial, como quando do homem se faz o cadaver. --- A mudança \emph{substancial} pode dar-se --- ou somente quanto á \emph{forma substancial}, permanecendo a materia, e assim são todas as mudanças naturaes, --- ou quanto á forma e á materia, de modo que \emph{toda} a substancia se mude noutra, e temos a \emph{transubstanciação}. Exemplo d'esta total mudança de uma substancia noutra não se encontra na philosophia, mas unicamente na Theologia revelada, a qual nos ensina que, no S. Sacrificio da Missa, toda a substancia do pão se muda no Corpo adoravel e toda a substancia do vinho se converte no Sangue precioso de N. S. Jesus Cristo. --- A mudança \emph{accidental} chama-se \emph{alteração}, e é a passagem do sujeito para um ser accidental, como quando um homem de ignorante se torna sabio.}
\pg{357} \emph{a) O ente, que é acto puro, é immutavel.} --- O ente, que é \emph{acto puro}, não contem em si nenhuma potencialidade. Ora o ente, que não contem em si nenhuma potencialidade, é immutavel; porque toda a mudança importa uma passagem da potencia para o acto. Logo o ente, que é \emph{acto puro}, é \emph{immutavel}.

\emph{b) O ente, que é immutavel, é acto puro.} --- O ente, que é \emph{immutavel}, não contem em si nenhuma potencialidade. Ora o ente, que não contem em si nenhuma potencialidade, é acto puro. Logo o ente, que é \emph{immutavel}, é \emph{acto puro} (¹).

\textbf{86. O ente mutavel não pode ser acto puro.} --- O ente \emph{mutavel} em tanto pode estar sujeito a mudanças, em quanto contem em si alguma potencialidade. Ora o que contem em si alguma potencialidade não é acto puro. Logo o ente \emph{mutavel} não pode ser \emph{acto puro}.\footnote{O ente que se muda, muda-se emquanto perde alguma coisa e adquire outra. Ora esta perda e esta acquisição não seria possivel, se não se desse uma \emph{composição} do sujeito que se muda, --- \emph{antes} da mudança, com a coisa que se perde, --- \emph{depois} da mudança, com a coisa que se adquire. Diz S. Thomaz: «Omne quod movetur, quantum ad aliquid manet, et quantum ad aliquid transit; sicut quod movetur de albedine in nigredinem, manet secundum substantiam. Et sic in omni quod movetur, attenditur aliqua compositio» (\emph{Sum}. \emph{Th}., q. I, q. 9, a. 1). Se \emph{o} ente mutavel é necessariamente \emph{composto}, não é \emph{acto puro}.}

\textbf{87. O ente immutavel é necessario, infinito, simplicissimo; o ente necessario, infinito, simplicissimo, é immutavel.}

\emph{a) O ente immutavel é necessario, infinito, simplicissimo.} --- O ente immutavel é acto puro. Ora o ente, que é acto puro, é necessario, infinito, simplicissimo. Logo o ente \emph{immutavel} é \emph{necessario}, \emph{infinito}, \emph{simplicissimo}.

\emph{b) O ente necessario, infinito, simplicissimo, é immutavel.} ---\footnote{Ninguem se admire se a demonstração d' estas proposições se basêa sempre neste criterio --- se um ente tem, ou não, \emph{potencialidade}. A essencia de tudo o que existe ou é \emph{acto} ou é composta de \emph{potencia} e de \emph{acto}. E, como a essencia é a raiz de todas as propriedades, é claro que a ella é preciso recorrer, como ao verdadeiro e intrinseco criterio, para conhecermos o que convem a um ente. --- \emph{O} ente, que é \emph{acto puro}, exclue toda a potencialidade, e por isso não é capaz nem de perder nem de adquirir alguma coisa, é \emph{immutavel}. --- O ente, que é \emph{immutavel}, não é capaz de adquirir ou de perder coisa alguma, por isso não tem potencialidade, é \emph{acto puro}. --- Diz S. Thomaz: «Primum ens oportet esse purum actum absque permixtione alicujus potentiae, eo quod potentia simpliciter est posterior actu. Omne autem quod quocumque modo mutatur, est aliquo modo in potentia. Ex quo patet quod impossibile est \textsc{Deum} aliquo modo mutari» (\emph{Sum}. \emph{Th}., q. I, q. 9, a. 1).}\pg{358} O ente necessario, infinito, semplicissimo, é só acto, sem nenhuma sombra de potencialidade. Ora o que é só acto, sem sombra de potencialidade, é immutavel. Logo o ente \emph{necessario}, \emph{infinito}, \emph{simplicissimo}, \emph{é immutavel} \footnote{O ente \emph{immutavel} --- é \emph{necessario}, porque, se fosse \emph{contingente}, receberia de outro ente a existencia, e pelo facto de passar da potencia para o acto, já não seria \emph{immutavel}. --- é \emph{infinito}, porque, se fosse \emph{finito}, estaria em potencia para uma perfeição ulterior, e deixaria de ser \emph{immutavel}. --- é \emph{simplicissimo}, porque, se fosse \emph{composto}, tambem deixaria de ser \emph{immutavel}. --- O ente \emph{necessario}, \emph{infinito}, \emph{simplicissimo}, é \emph{immutavel}; porque, --- se o ente \emph{necessario} fosse \emph{mutavel}, ou adquiriria uma perfeição que não tinha, ou perderia uma perfeição que tinha, e assim não seria \emph{necessario}, mas \emph{contingente} em relação á perfeição adquirida ou perdida, --- se o ente \emph{infinito} fosse \emph{mutavel}, tambem ou perderia uma perfeição, e não seria \emph{infinito}, mas \emph{finito}, --- se o ente \emph{simplicissimo} fosse \emph{mutavel}, passaria da potencia para o acto, e já não seria \emph{simplicissimo}, mas \emph{composto} (Cf. \emph{Sum}. \emph{Th}., l. c.). Como se vê, \emph{acto puro}, \emph{ente necessario}, \emph{ente infinito}, \emph{ente simplicissimo}, \emph{ente immutavel}, são attributos que exprimem uma e a mesma coisa, e por isso a um ente, ao qual convenha um d'esses attributos, convêm todos, e de um d'elles pode legitimamente deduzir-se outro.}. §8. \textbf{O ente mutavel é contingente, finito, composto; o ente contingente, finito, composto, é mutavel.}

a) \emph{O ente mutavel é contingente, finito, composto.} --- Ente mutavel é o que está sujeito a mudanças, e por isso não é acto puro. Ora o que não é acto puro é contingente, finito, composto. Logo o ente \emph{mutavel é contingente, finito, composto}.

b) \emph{O ente contingente, finito, composto, é mutavel.} --- Todo o ente contingente, finito, composto, não é acto puro, mas tem em si alguma potencialidade. Ora o que tem potencialidade é mutavel. Logo o ente \emph{contingente, finito, composto, é mutavel} \footnote{O ente \emph{mutavel} --- é \emph{contingente}, porque, na mudança, adquire ou perde uma perfeição, e assim é \emph{contingente} em relação á perfeição adquirida ou perdida, --- é \emph{finito}, porque, se, na mudança, adquire ou perde uma perfeição, não possue todas as perfeições, e por isso é \emph{finito}, --- é \emph{composto}, porque, quando \emph{muda}, passa da potencia para o acto, e por isso consta d'esses dois elementos. --- O ente \emph{contingente, finito, composto, é mutavel}. É mutavel o ente \emph{contingente}; porque, sendo \emph{contingente}, passou e pode passar da potencia para o acto, e assim mudar-se. É mutavel o ente \emph{finito}; porque, sendo \emph{finito}, pode adquirir uma perfeição ulterior, e assim passar da potencia para o acto. É mutavel o ente \emph{composto}; porque, sendo composto de potencia e de acto, pode passar de um estado para outro. (\emph{Sum}. \emph{Th}., l. c., a. 2).}.

§9. \textbf{O ente mutavel suppõe o ente immutavel.} --- Na ordem \emph{ontologica}, ou da \emph{actualidade}, o ente \emph{actual} deve preceder o ente \emph{possivel}; porque este só pode adquirir o acto pela acção de um\pg{359} ente, que já o possue. Ora o ente \emph{mutavel} é \emph{possivel}, ou está em potencia em relação ao acto, para o qual se muda. Logo o. ente \emph{mutavel} deve ser reduzido ao acto por outro ente. --- Este outro ente ou é immutavel ou é mutavel: Se é immutavel. está provada a conclusão. Se é mutavel, é \emph{possivel} e por isso suppõe um outro ente, e este um outro, até se encontrar um immutavel; pois, se não existisse o ente immutavel, não poderia existir nenhum ente mutavel. Logo o ente \emph{mutavel} suppõe o ente immutavel (1).

90. Corollarios. --- Dos principios expostos resulta:

I. Na ordem \emph{ontologica}, ou da \emph{realidade}, o immutavel precede o \emph{mutavel}; mas, na ordem \emph{logica}, ou do \emph{conhecimento}, o \emph{mutavel} precede o \emph{immutavel}, porque é do conhecimento d'aquelle que nos elevamos ao conhecimento d'este.

II. Dada a existencia de um só ente \emph{mutavel}, deve necessariamente admittir-se a existencia de um ente immutavel\footnote{Vemos no mundo entes sujeitos a mudanças, que do não-ser passaram para o ser. Esta passagem não pode ter sido realisada pelos proprios entes, que, estando em potencia, isto é, sendo possiveis mas não extistentes, não podiam operar; pois á operação é uma consequencia do ser. Logo essa passagem deve ter a sua causa proporcionada na força de um ente exterior, e, em ultima analyse, absolutamente immutavel: porque, se não existisse o \emph{immutavel}, não poderia existir, como effectivamente existe, o \emph{mutavel}. Diz S. Thomaz: «Videmus aliquid esse in mundo, quod exit de potentia in actum. Non autem educit se de potentia in actum; quia quod est (in) potentia, nondum est, unde nec agere potest. Ergo oportet esse aliquid aliud prius, quo educatur de potentia in actum. Et iterum, si hoc est exiens de potentia in actum, oportet ante hoc aliquid aliud poni, quo reducatur in actum. Hoc autem in infinitum procedere non potest. Ergo oportet devenire ad aliquid, quod est tantum actu, et nullo modo in potentia; et hoc dicimus Deum» (\emph{C. Gent.}, I. 16).}\footnote{Diz S. Fulgencio: «Firmissime tene et nullatenus dubites omnem creaturam naturaliter mutabilem a Deo incommutabili factam» (\emph{De fide ad Petrum}, c. 20). --- Recommendamos aos leitores o estudo da doutrina, exposta neste capitulo segundo: pois constitue ella a base dos argumentos, que provam a existencia de Deus.}
